% BHU PYQ -- Manually transcribed & error-corrected
% Paper: MTB-602 : Linear Algebra
% Exam: B.A./B.Sc. (Hons) Semester VI Examination, 2013-14

\documentclass[11pt,a4paper]{article}
\usepackage[a4paper,top=2.5cm,bottom=2.5cm,left=2.8cm,right=2.8cm,headheight=14pt]{geometry}
\usepackage{amsmath,amssymb}
\usepackage[T1]{fontenc}
\usepackage[utf8]{inputenc}
\usepackage{lmodern}
\usepackage{microtype}
\usepackage{fancyhdr}
\usepackage{enumitem}
\usepackage{hyperref}
\usepackage{lastpage}
\usepackage{parskip}

\hypersetup{colorlinks=true,urlcolor=black,linkcolor=black,
  pdftitle={MTB-602: Linear Algebra -- BHU B.A./B.Sc. Sem VI 2013-14}}
\pagestyle{fancy}\fancyhf{}
\renewcommand{\headrulewidth}{0.4pt}\renewcommand{\footrulewidth}{0.4pt}
\fancyhead[L]{\small   \textit{Banaras Hindu University}}
\fancyhead[R]{\small   \textit{MTB-602: Linear Algebra}}
\fancyfoot[C]{\small Page  \thepage\ of \pageref{LastPage}}


\newlist{parts}{enumerate}{2}
\setlist[parts,1]{label=(\alph*),leftmargin=*,itemsep=5pt,topsep=3pt}
\setlist[parts,2]{label=(\roman*),leftmargin=*,itemsep=3pt,topsep=2pt}

\newcommand{\pts}[1]{\hfill   \textit{[#1]}}


\begin{document}

\begin{center}
  {\large\bfseries BANARAS HINDU UNIVERSITY}\\[2pt]
  {
\normalsize B.A./B.Sc. (Hons) --- Semester VI Examination, 2013-14}\\[6pt]
  \rule{0.8\linewidth}{0.5pt}\\[6pt]
  {\large\bfseries Paper No. MTB-602: Linear Algebra}\\[4pt]
  {
\normalsize Time: 3 Hours \hfill Full Marks: 70}
\end{center}
\rule{\linewidth}{0.4pt}
\smallskip



\noindent   \textit{Instructions: Answer   \textbf{five} questions including Question~1 which is compulsory. Figures in right-hand margin indicate marks.}
\bigskip




\noindent   \textbf{Question 1.}   \textit{(Compulsory --- 2 marks each.)}

\begin{parts}
  \item For a vector space $V$ over a field $F$ and $x, y \in V$, $a \in F$, show that:
  
\begin{parts}
    \item $a(x-y) = ax - ay$
    \item $ax = 0 \implies a = 0$ or $x = 0$
  \end{parts}
  \item Show that every superset of a linearly dependent set of vectors is linearly dependent.
  \item Show that the vectors $(1, 2, 1)$, $(2, 1, 0)$, and $(1, -1, 2)$ form a basis of the vector space $\mathbb{R}^3$ over $\mathbb{R}$.
  \item Let $T$ be a linear transformation on a vector space $V$ such that $T^2 - T + I = 0$. Show that $T$ is invertible.
  \item Show that similar matrices have the same determinant.
  \item If $v$ is a characteristic vector of a linear operator $T$, show that $v$ cannot correspond to more than one characteristic root of $T$.
  \item If $x$ and $y$ are vectors in a unitary space $V$, then prove that $(x, y) =  \text{Re}(x, y) + i  \text{Re}(x, iy)$.
\end{parts}

\medskip\rule{\linewidth}{0.2pt}

\medskip


\noindent   \textbf{Question 2.}

\begin{parts}
  \item Show that a non-empty subset $W$ of a vector space $V$ is a subspace of $V$ if and only if for all $x, y \in W$ and $a \in F$, $ax + y \in W$. Hence determine whether the set of ordered triads $(x, y, z)$ of real numbers with $z > 0$ is a subspace of $\mathbb{R}^3(\mathbb{R})$ or not. \pts{7}
  \item Show that every linearly independent subset of a finitely generated vector space $V(F)$ is either a basis of $V$ or can be extended to form a basis of $V$. \pts{7}
\end{parts}

\medskip\rule{\linewidth}{0.2pt}

\medskip


\noindent   \textbf{Question 3.}

\begin{parts}
  \item If $S$ and $T$ are subsets of a vector space $V(F)$, then prove that:
  
\begin{parts}
    \item $L(S \cup T) = L(S) + L(T)$
    \item $S$ is a subspace of $V$ if and only if $L(S) = S$.
  \end{parts} \pts{7}
  \item State and prove the rank-nullity theorem for linear transformations. \pts{7}
\end{parts}

\medskip\rule{\linewidth}{0.2pt}

\medskip


\noindent   \textbf{Question 4.}

\begin{parts}
  \item Let $T: V  o V$ be a linear transformation. Prove that the following conditions are equivalent:
  
\begin{parts}
    \item $R(T) \cap N(T) = \{0\}$
    \item $T(T(x)) = 0 \implies T(x) = 0$
  \end{parts} \pts{6}
  \item Find the matrix of the linear transformation $T: \mathbb{R}^3  o \mathbb{R}^3$ defined by:
  \[
  T(a, b, c) = (2b + c, a - 4b, 3a)
  \]
  with respect to the ordered basis $B = \{(1, 1, 1), (1, 1, 0), (1, 0, 0)\}$. \pts{8}
\end{parts}

\medskip\rule{\linewidth}{0.2pt}

\medskip


\noindent   \textbf{Question 5.}

\begin{parts}
  \item Let $U$ and $V$ be finite dimensional vector spaces over a field $F$ of dimensions $n$ and $m$ respectively. Prove that the vector space $L(U, V)$ of all linear transformations from $U$ to $V$ is finite dimensional and has dimension $mn$. \pts{7}
  \item Let $V$ be an $n$-dimensional vector space over a field $F$ and let $B = \{x_1, x_2, \dots, x_n\}$ be an ordered basis for $V$. If $\{a_1, a_2, \dots, a_n\}$ is any ordered set of $n$ scalars from $F$, show that there exists a unique linear functional $f$ on $V$ such that $f(x_i) = a_i$ for $i = 1, 2, \dots, n$. \pts{7}
\end{parts}

\medskip\rule{\linewidth}{0.2pt}

\medskip


\noindent   \textbf{Question 6.}

\begin{parts}
  \item Diagonalize the following matrix:
  \[
  A = 
\begin{pmatrix} 9 & -4 & -4 \\ -8 & 3 & 4 \\ -16 & 8 & 7 \end{pmatrix}
  \] \pts{8}
  \item Define an inner product space. Let $x = (x_1, x_2, \dots, x_n)$ and $y = (y_1, y_2, \dots, y_n)$ in $\mathbb{C}^n$. Define $(x, y) = x_1 ar{y}_1 + x_2 ar{y}_2 + \dots + x_n ar{y}_n$. Show that $\mathbb{C}^n$ is an inner product space. \pts{6}
\end{parts}

\medskip\rule{\linewidth}{0.2pt}

\medskip


\noindent   \textbf{Question 7.}

\begin{parts}
  \item State and prove the Cauchy-Schwarz inequality. Find when this inequality becomes an equality and supply a proof. \pts{7}
  \item Define an orthonormal set in an inner product space. Show that:
  
\begin{parts}
    \item Every finite-dimensional inner product space $V 
eq \{0\}$ possesses an orthonormal basis.
    \item If $S = \{x_1, x_2, \dots, x_n\}$ is an orthonormal set of vectors in an inner product space $V$ and if a vector $y$ is in the linear span of $S$, then:
    \[
    y = \sum_{k=1}^n (y, x_k) x_k
    \]
  \end{parts} \pts{7}
\end{parts}

\begin{center}
   \textbf{OR}
\end{center}



\noindent   \textbf{Question 7 [Alternate].} Apply the Gram-Schmidt process to the vectors $(1, 0, 1)$, $(1, 3, 1)$, and $(3, 2, 1)$ in $\mathbb{R}^3$ to obtain an orthonormal basis of the inner product space $\mathbb{R}^3$ over $\mathbb{R}$ with respect to the standard inner product. \pts{7}

\vfill
\rule{\linewidth}{0.4pt}

\begin{center}\small
  \href{https://bhu-exam.vercel.app/}{Click here to visit our website}\\[4pt]
  \href{https://me-aryan.vercel.app/}{Click here to visit our contributor}\\[4pt]
  \href{https://quashan-me.vercel.app/}{Click here to visit our contributor}
\end{center}

\end{document}
